\section{What actual lease disclosures identify}
\label{sec:henkel-lease-measurements}

The preceding allocation example starts with named rights, supplied contract
findings and a hypothetical opening book. A company report usually supplies
less. Before using it to select an opening, the simulator must distinguish what
was reported from the contract state needed to calculate the next payment or
modification. Henkel's retained 2025 annual report makes that distinction
concrete. It gives several useful lease measurements, but those measurements
neither enumerate the individual rights nor allocate their carrying amounts.

The following amounts are in millions of euros for the consolidated Henkel
Group. The report describes its basis as IFRS Accounting Standards as adopted
in the European Union. Its 2024 comparative and 2025 current figures belong to
the same retained 2025 document; their appearance together supplies no evidence
that these exact bytes were public in 2024. The basis note also says that
individual figures are rounded and addition can differ from the indicated
total. It supplies no rounding tie rule.\footnote{Henkel, \emph{Annual Report
2025}, printed p.~351, physical PDF p.~353. The retained file is
\path{data/financial-report-benchmark/documents/henkel/henkel-2025-annual/original.pdf}.
The source audit records its SHA-256, exact text spans and table-word geometry
under \path{data/accounting-master/henkel-lease-source-audit-20261024}. All page
references in this section are printed page numbers; physical PDF pages are
two higher.}

At 31 December 2025, Henkel reports right-of-use assets of 594, lease
liabilities of 666 and undiscounted contractual lease cash flows of 721.
These amounts answer different questions. The first measures the remaining
carrying value of the rights to use assets; the second measures the recognised
liability; the third groups future undiscounted payments by maturity. The
liability consists of 557 noncurrent and 109 current. The maturity bands are
114 due within one year, 351 due between one and five years, and 255 due after
five years. Their displayed sum is 720, one below the reported total of
721.\footnote{Henkel, \emph{Annual Report 2025}, pp.~372--374, 397 and 423--424.
The maturity table is explicitly based on undiscounted cash flows.}

Different measurement bases explain why 594, 666 and 721 should not be forced
to agree. Rounding creates a separate issue even within a single basis.
The reported 2025 net right-of-use amounts for land and buildings, plant and
machinery, and factory and office equipment are 530, 16 and 47. Those three
numerical entries sum to 593; the total is 594. The construction category is
printed as a dash. The inventory keeps that dash distinct from zero and records
the difference between the total and the three numerical entries without
inventing a fourth amount. Similarly, the 2024 gross right-of-use amount is
1,173 and accumulated depreciation and impairment is 543, while net carrying
amount is 631. Treating every displayed integer as an exact underlying amount
would demand both $1{,}173-543=630$ and a net amount of 631. That exact-value
interpretation is wrong; it does not establish an error in the rounded report.

\subsection{Read the liability bridge with its signs intact}

The cash-flow reconciliation helps distinguish a payment from the expense
recognised in the same year. Henkel presents liabilities with a negative sign
in this reconciliation. A payment reducing the liability therefore enters
positively; an addition or accrued interest increasing it enters negatively.
For 2025, the displayed opening liability is $-713$. Cash changes are $158$,
comprising interest paid of 25 and redemption of 133. Interest expense is a
separate movement of $-25$, additions are $-139$, acquisitions or disposals of
subsidiaries contribute 1, foreign exchange contributes 50, and sundry changes
contribute 2. The fair-value column displays \texttt{-0}.
Using the displayed numerical value of that last entry gives
\begin{equation}
\label{eq:henkel-displayed-liability-bridge}
 -713+158-25-139+1+50+0+2=-666.
\end{equation}
This is exact arithmetic on the displayed entries, in EUR millions, for the
calendar year 2025. The right side agrees with the negative of the reported
closing liability. The sign of the printed zero remains in the source record
even though its displayed rational value is zero.\footnote{Henkel,
\emph{Annual Report 2025}, pp.~450--451. The cash-flow statement on p.~347
reports redemption as an outflow of $-133$ in 2025 and $-146$ in 2024.}

The subtotal 158 enters equation~\eqref{eq:henkel-displayed-liability-bridge}
once. Adding 25 and 133 again would count the same cash twice. The interest
expense of $-25$ has a different role: it increases the liability that the
interest payment reduces. The 2024 reconciliation follows the same sign
convention, but the sum of its displayed opening and movements is $-712$,
whereas its displayed closing balance is $-713$. The retained residual is
$-1$. The rounding disclosure makes this a reason to examine precision, not
authority to insert a missing transaction.

Nor does the bridge explain every lease-related cash flow. Henkel reports
total cash paid for leases of 174 in 2025, exceeding the recognised-liability
cash subtotal of 158 by 16. Short-term lease expense, low-value lease expense
and excluded variable-payment expense are 11, 5 and 1, respectively. Their
sum is 17, and the report does not identify them as the allocation of that cash
difference. Expenses and cash flows can differ in both timing and scope. The
source inventory retains those measurements separately.

The same care is necessary for the asset movement. Opening net right-of-use
assets of 631, additions of 139, and depreciation and impairment of 132 give
$631+139-132=638$, compared with a reported closing amount of 594. Those three
movements are not a complete asset reconciliation. The difference of $-44$
does not identify a particular disposal, exchange effect or allocation to a
right. It remains a diagnostic of the incomplete comparison.

\subsection{Ask a conditional question about unrounded measurements}

The useful next question is whether one set of unrounded measurements could
satisfy selected relations while producing the displayed figures. Answering
it requires an explicit rounding hypothesis and explicit premises about which
underlying relations are exact. The report's general rounding statement does
not choose these premises for us.

For illustration, suppose that each selected amount is rounded to the nearest
million euros and exact half ties are rounded away from zero. Let $d$ denote
a displayed integer and $x$ the corresponding unrounded amount, both measured
in EUR millions. Under this proposed rule the inverse interval is
\begin{equation}
\label{eq:henkel-proposed-rounding}
 x\in I(d)=
 \begin{cases}
 [d-\tfrac12,d+\tfrac12),&d>0,\\
 (d-\tfrac12,d+\tfrac12],&d<0,\\
 (-\tfrac12,\tfrac12),&d=0.
 \end{cases}
\end{equation}
At the positive lower endpoint, a half tie rounds up to $d$; at the positive
upper endpoint it rounds to $d+1$. For a negative number the included endpoint
is reversed. Neither half endpoint rounds to zero. Thus the zero case is open
at both ends. Treating the printed \texttt{-0} numerically as zero imposes no
additional sign restriction: the source has not established what sign rule,
if any, its printed zero encodes.

This hypothesis can reconcile the 2024 gross, accumulated and net amounts
without changing their displayed values. For example, underlying amounts
1,173.25, 542.50 and 630.75 satisfy
$1{,}173.25-542.50=630.75$ and round to 1,173, 543 and 631 under the proposed
rule. The example is a conditional mathematical assignment, not a recovery
of Henkel's actual amounts. A different rounding convention would require
different endpoint rules.

Compatibility must also hold jointly. Let $L_{25}$ be the positive closing
liability, $L^{\mathrm{nc}}_{25}$ and $L^{\mathrm{cur}}_{25}$ its noncurrent and
current portions, $B_{25}$ its signed reconciliation closing, and $O_{25}$
the opening of the 2025 reconciliation. Let $B_{24}$
denote the signed closing at 31 December 2024. Selected underlying premises
include
\begin{equation}
\label{eq:henkel-shared-measurements}
 L_{25}=L^{\mathrm{nc}}_{25}+L^{\mathrm{cur}}_{25},
 \qquad B_{25}=-L_{25},
 \qquad O_{25}=B_{24}.
\end{equation}
The same $B_{25}$ appears as the closing amount in the 2025 stock-flow
relation, and the same $O_{25}$ appears there as its opening. Separate checks
could choose incompatible values for these shared amounts. One joint system
requires a common assignment. The opening is dated 31 December 2024, the
movements cover $[2025\text{-}01\text{-}01,2026\text{-}01\text{-}01)$, and
the closing is dated 31 December 2025; the opening symbol is indexed by the
reconciliation year, not its measurement year.

The checked reference combines nineteen such relations over fifty-four
selected measurements from the two years. With the proposed rounding and
identity premises, it finds an assignment satisfying every interval and
equation, checked with exact rational arithmetic. This establishes conditional
consistency of the selected measurements. It neither chooses the actual
unrounded amounts nor makes the rounded intervals independent. The 2024
displayed-as-exact counterexample fails because its three exact values cannot
satisfy the gross-to-net relation; that failed hypothesis is distinct from a
contradiction in the underlying accounts.

\subsection{Preserve the disclosures before admitting a lease book}

\path{HenkelLeaseMeasurements/1} binds this finite example to its retained
source. Its closed request contains \path{profile}, \path{inventory_sha256},
\path{selected_relations}, \path{rounding_hypothesis} and
\path{identity_hypotheses}. The hash selects the exact inventory. Relation IDs
must be unique members of the supported registry; the identity map contains
exactly those IDs. Each premise is absent or is the explicitly proposed
premise for that relation. Callers cannot replace the amounts, units, dates,
company or equations through this interface.

For each presentation year, the registry supports current plus noncurrent
liability; gross less accumulated right-of-use amounts; depreciation categories;
undiscounted maturity bands; cash subtotal components; the signed liability
stock-flow; the closing-liability sign comparison; the redemption sign
comparison; and the maturity table's separate carrying-value comparison.
Opening-to-closing continuity between the two years supplies the nineteenth
relation. Category sums involving unquantified dashes, the incomplete asset
movement, attribution of the cash difference to expenses, and equality between
discounted and undiscounted amounts are excluded.

Every response retains all 104 numerical-table and narrative entries, including
eight dash cells, the twelve qualitative observations and fourteen unresolved
state dependencies. The selected measurements carry their original source
references, literals, dates, units and consolidation context. A stock-flow
relation retains the roles of opening, period movements and closing; these
roles are validated before lowering to the existing rational constraint engine.
Shared observations keep one identity throughout the selected system.

By default, the adapter supplies no numeric inverse intervals. Missing rounding
or identity premises leave the assessment \path{unresolved}. An empty relation
selection reports \path{no_asserted_system} while retaining the source record.
A verified assignment gives \path{conditionally_feasible}; a checked
contradiction concerns only the declared conditional system. Resource exhaustion
returns \path{unknown}, and invalid source bindings or declarations return
\path{invalid}. The reference permits at most one conjunction cell, with an
explicit bounded solver budget; zero cell budget supplies no verified witness.
Saved results are rebuilt from the original request. Fixed fields must agree,
and saved witnesses must satisfy the exact intervals and equations; a different
valid witness or elapsed runtime need not have identical bytes.

An actual lease book still needs individual lease and right identities, complete
membership, legal lessees and owner books, unrounded allocations, contract and
functional currencies, rates, commencement and measurement dates, payment
schedules, modifications, depreciation and applicable findings. Henkel's
description of country-, duration- and credit-specific discount rates does not
supply those numeric inputs. Its description of optional lease periods does
not supply contract-specific exercise findings, and its qualitative
mid-triple-digit-million-euro payment description remains qualitative.
Short-term and low-value recognition exemptions require separate treatment.
Initial right-of-use measurement can also include payments, incentives and
other adjustments beyond the liability; equality between the two is not a
general allocation rule.\footnote{Henkel, \emph{Annual Report 2025},
pp.~373--374. The retained disclosure describes these policies; the present
example does not independently certify their application to each contract.}

The source adapter therefore leaves executable opening state absent and source,
accounting and native-state admission false. It preserves useful measurements
while keeping their missing contract explanation visible. The synthetic USD
lease methods retain their original company, period and authority scope; they
cannot become Henkel methods by changing a name or currency label. The full
requests, relations and exact checked results are retained under
\path{data/accounting-master/henkel-lease-measurements-20261024}, with the
reference in \path{scripts/henkel_lease_measurements_v1.py}. Connecting these
measurements to actual contracts remains the next substantive reconstruction
problem.
